\chapter{Evaluación y resultados}
\label{chap:evaluation_results}

\section{Diseño de experimentos y métricas de evaluación}
\label{sec:experiment_design}
% TODO: Describir el diseño experimental completo: definición de hipótesis, variables dependientes e independientes, métricas de evaluación seleccionadas (precisión OCR, latencia, usabilidad), y metodología estadística para análisis de resultados.

\section{Resultados de precisión, latencia y eficiencia energética}
\label{sec:performance_results}
% TODO: Presentar resultados cuantitativos del modelo OCR: precisión por tipo de dígito, matriz de confusión, tiempos de inferencia en diferentes dispositivos, consumo energético, y comparativa con modelos baseline.

\section{Comparativa con procesos tradicionales de lectura}
\label{sec:traditional_comparison}
% TODO: Analizar comparativamente el sistema desarrollado versus el método manual: tiempo de lectura, tasa de errores, costo operativo, y impacto en la eficiencia de las JAAPS.

\section{Validación en campo con usuarios reales}
\label{sec:field_validation}
% TODO: Documentar resultados de la validación en JAAPS reales: perfil de usuarios participantes, escenarios de uso, métricas de usabilidad (SUS), aceptación tecnológica, y feedback cualitativo de los operadores.

\section{Discusión de resultados y cumplimiento de objetivos}
\label{sec:results_discussion}
% TODO: Analizar críticamente los resultados obtenidos, evaluar el cumplimiento de objetivos específicos planteados, identificar factores que influyeron en el rendimiento, y discutir limitaciones y fortalezas del sistema desarrollado.